%% ma: L11-B02-0005
%% lop: 11
%% chuong: 1
%% bai: 2
%% dang: 11.2.8
%% loai: TL
%% muc: VD
%% dap_an: "-"
%% nguon: "de-cuong-giua-ky-1-toan-11-nam-2026-2027-truong-thpt-hoang-van-thu-ha-noi-2-8.pdf (M08, Đề cương giữa kì 1 Toán 11 THPT Hoàng Văn Thụ, Hà Nội), Phần 4 Bài 1 ý 2"
%% kien_thuc: "Bài 1 (giá trị lượng giác của góc lượng giác), Bài 2 (công thức lượng giác)"
%% trang_thai: cho-duyet
%% ly_do: "Dạng mới đề xuất 11.2.8"
%% ghi_chu: "Đề không cho góc phần tư của α nên sin 4α có hai giá trị ±3√7/8."
\begin{ex}
Cho $\sin\alpha+\cos\alpha=\dfrac12$. Tính $\sin2\alpha$; $\sin4\alpha$.
\loigiai{$(\sin\alpha+\cos\alpha)^2=1+\sin2\alpha=\dfrac14\Rightarrow\sin2\alpha=-\dfrac34$.

$\cos^22\alpha=1-\dfrac9{16}=\dfrac7{16}\Rightarrow\cos2\alpha=\pm\dfrac{\sqrt7}4$ (cả hai dấu đều xảy ra vì $\cos2\alpha=(\cos\alpha-\sin\alpha)\cdot\dfrac12$ với $\cos\alpha-\sin\alpha=\pm\dfrac{\sqrt7}2$).

$\sin4\alpha=2\sin2\alpha\cos2\alpha=\mp\dfrac{3\sqrt7}{8}$, tức là $\sin4\alpha=\dfrac{3\sqrt7}8$ hoặc $\sin4\alpha=-\dfrac{3\sqrt7}8$.}
\end{ex}
